\documentclass[11pt,letterpaper]{article}
\usepackage[T1]{fontenc}
\usepackage[utf8]{inputenc}
\usepackage{amsmath,amsthm,mathtools}
\usepackage{newtxtext,newtxmath}
\usepackage{bm}
\usepackage[margin=1in,headheight=14pt]{geometry}
\usepackage{microtype}
\usepackage{booktabs,array,longtable}
\usepackage[dvipsnames]{xcolor}
\usepackage{enumitem}
\usepackage{listings}
\usepackage[most]{tcolorbox}
\usepackage{fancyhdr}
\usepackage{hyperref}
\usepackage[nameinlink,noabbrev]{cleveref}
\definecolor{ink}{HTML}{18334D}
\definecolor{accent}{HTML}{276B67}
\definecolor{pale}{HTML}{F1F5F7}
\hypersetup{colorlinks=true,linkcolor=ink,citecolor=accent,urlcolor=accent,
 pdftitle={A Twelve-Value Proof of the A181280 Formula},
 pdfauthor={ChatGPT},pdfsubject={Exact enumeration, binary Gram matrices, Fourier recurrences}}
\pagestyle{fancy}
\fancyhf{}
\fancyhead[L]{\small\color{ink}A twelve-value proof of A181280}
\fancyhead[R]{\small\color{ink}Research report}
\fancyfoot[C]{\thepage}
\renewcommand{\headrulewidth}{0.3pt}
\setlength{\parskip}{4pt}
\setlength{\parindent}{1.2em}
\setlist{itemsep=2pt,topsep=4pt}
\newtheorem{theorem}{Theorem}[section]
\newtheorem{lemma}[theorem]{Lemma}
\newtheorem{proposition}[theorem]{Proposition}
\newtheorem{corollary}[theorem]{Corollary}
\theoremstyle{definition}
\newtheorem{definition}[theorem]{Definition}
\newtheorem{remark}[theorem]{Remark}
\newtheorem{certificate}[theorem]{Finite certificate}
\newcommand{\F}{\mathbb F_2}
\newcommand{\Q}{\mathbb Q}
\newcommand{\Sym}{\operatorname{Sym}}
\newcommand{\rank}{\operatorname{rank}}
\newcommand{\rad}{\operatorname{rad}}
\newcommand{\Astar}{\mathcal A_*}
\newcommand{\1}{\mathbf 1}
\newcommand{\e}{\mathbf e}
\newcommand{\ip}[2]{\langle #1,#2\rangle}
\newcommand{\lex}{<_{\rm lex}}
\newcommand{\revlex}{>_{\rm lex}}
\newcommand{\good}{\mathsf{good}}
\newcommand{\allstate}{\{1,\ldots,m-1\}}
\newtcolorbox{researchbox}[1][]{colback=pale,colframe=accent,boxrule=0.5pt,
 arc=1mm,left=3mm,right=3mm,top=2mm,bottom=2mm,#1}
\lstset{language=Python,basicstyle=\ttfamily\footnotesize,
 keywordstyle=\color{ink}\bfseries,commentstyle=\color{accent},
 stringstyle=\color{Brown},columns=fullflexible,keepspaces=true,
 showstringspaces=false,breaklines=true,frame=single,rulecolor=\color{black!18},
 xleftmargin=2mm,xrightmargin=2mm,aboveskip=8pt,belowskip=8pt}
\allowdisplaybreaks[1]

\title{\vspace{-1cm}\color{ink}\bfseries
 A Twelve-Value Proof of the A181280 Formula\\[6pt]
 \large Fourier spectra and universal recurrences for ordered binary Gram matrices}
\author{ChatGPT\\\small AI-assisted mathematical research report}
\date{19 September 2026}

\begin{document}
\maketitle
\begin{abstract}
OEIS A181280 counts binary matrices with four lexicographically increasing rows
whose mod-2 Gram matrix has lexicographically decreasing rows. Kauers and
Koutschan proposed an explicit six-term exponential-polynomial formula in 2023.
We give a computer-assisted proof of that formula: a structural argument first
proves a universal order-12 recurrence for the tail, after which twelve exact
counts settle the identity for every $n\ge4$. The finite arithmetic is supplied
in two independent implementations, with a further Fourier audit and direct
brute-force checks.

The structural argument works for any fixed number $m$ of rows and any set of
allowed symmetric Gram matrices. It gives an explicit eventual recurrence of
order $m+2\lfloor m^2/4\rfloor$, with characteristic roots restricted to signed
powers of two and explicitly bounded multiplicities. For four rows this
universal order is sharp. We also determine the rational generating function,
the minimal eventual recurrence of order 11 for A181280, the exact transient,
and quantitative equidistribution and asymptotic consequences.
\end{abstract}

\begin{researchbox}[title=\textbf{Status and attribution}]
The OEIS formula was still labeled conjectural when retrieved on 19 September
2026~\cite{oeis}. However, Kauers's December 2025 slides already prove the general
\emph{existence} of a constant-coefficient recurrence~\cite{slides}. That prior result is
credited here, not claimed as new. The slides display the explicit formula with
a question mark but also mark A181280 with a check in an overview table; they do
not supply an explicit formula certificate. Thus the exact prior-resolution
status is not completely settled by the available sources. This report provides
a self-contained, reproducible proof, but makes no claim of priority and has not
received independent peer review.
\end{researchbox}

\noindent\textbf{Keywords.} OEIS A181280; finite-state enumeration; binary quadratic
forms; Walsh sums; C-finite sequences; rational generating functions.

\newpage
{\small\setlength{\parskip}{0pt}\tableofcontents}
\newpage

\section{The selected problem and the result}
\label{sec:problem}

\subsection{The counting problem}
For a binary word, lexicographic order uses $0<1$ and reads from left to right.
All words being compared have the same length. Let $a(n)$ be the number of
matrices $M\in\F^{4\times n}$ such that
\begin{equation}\label{eq:conditions}
 M_{1,*}\lex M_{2,*}\lex M_{3,*}\lex M_{4,*},
 \qquad
 (MM^\top)_{1,*}\revlex (MM^\top)_{2,*}
 \revlex (MM^\top)_{3,*}\revlex (MM^\top)_{4,*}.
\end{equation}
The product is taken over $\F$, not over the integers. Thus its $(i,j)$ entry
is the parity of the ordinary dot product of rows $i$ and $j$.
For $n\ge1$ this is the sequence A181280, introduced in the OEIS by R.\ H.\ Hardin
\cite{oeis}. We extend it by $a(0)=0$: four empty rows are not strictly ordered.

Kauers and Koutschan state the proposed formula as Conjecture 20 in
Section 6.5 of their 2023 paper \cite{kk}. The following is that same formula,
rewritten to group equal exponential bases.

\begin{theorem}[The A181280 formula]\label{thm:main}
For every integer $n\ge4$,
\begin{align}
 a(n)={}&\frac{16^n}{512}
 +\frac{288n-3473}{147456}\,8^n
 -\frac{113}{49152}\,(-8)^n \notag\\
 &+\frac{6n^2-219n+820}{6144}\,4^n
 -\frac{13n-164}{6144}\,(-4)^n
 -\frac{3n+32}{288}\,2^n.
 \label{eq:formula}
\end{align}
Also $a(0)=a(1)=a(2)=a(3)=0$.
\end{theorem}

The proof occupies Sections \ref{sec:automaton}--\ref{sec:certificate}.
The only computer-assisted part essential to the theorem is an exact finite
enumeration through $n=15$. The reason this is sufficient is proved before the
counts are used. No guessed recurrence is assumed in that proof.

\subsection{What was already known, and what is established here}
Kauers's talk \emph{Guessing with little data}, listed in the Rutgers seminar
archive for 11 December 2025 \cite{seminar}, contains a transfer-matrix and
parity-filtering proof that these counting sequences are C-finite for arbitrary
fixed row count and arbitrary allowed Gram set \cite{slides}. Here
\emph{C-finite} means satisfying a linear recurrence with constant coefficients.
Its formula slide is numbered 30, and its general C-finiteness theorem is on
slide 39. Because the PDF contains overlay pages, these are PDF pages 104 and
137, respectively, counting pages from one.

The contribution developed in this report is more explicit than existence:
we locate the possible characteristic roots, bound their multiplicities using
a filtration with only $m$ levels, prove a quadratic-in-$m$ universal recurrence
order, and use that bound to finish a finite certificate for the displayed
formula. We then prove sharpness of the four-row universal bound and derive
several exact consequences. The literature check located the 2023 conjecture,
the live OEIS entry, and the 2025 slides, but not a separate complete proof of
the explicit formula. That search is not a proof that no such proof exists.

\subsection{The key reduction}
A direct automaton would have $8\cdot1024=8192$ states: eight row-order states
and 1024 possible symmetric $4\times4$ Gram matrices. The key fact is that its
counting sequences have a much smaller universal annihilator. With $E$ denoting
the forward shift, $(Ef)(n)=f(n+1)$, define
\begin{align}
 P_4(x)={}&(x-2)^2(x+2)(x-4)^3(x+4)^2\notag\\[-2pt]
 &\hspace{20mm}\cdot(x-8)^2(x+8)(x-16).
 \label{eq:P4}
\end{align}
We prove
\begin{equation}\label{eq:keyrec}
 P_4(E)a(n)=0\qquad(n\ge4).
\end{equation}
This monic polynomial has degree 12. The right side of \eqref{eq:formula}
satisfies the same recurrence. Therefore equality at $n=4,5,\ldots,15$ proves
equality at every later index by induction. Establishing \eqref{eq:keyrec}
from the definition, rather than inferring it from initial data, is what turns
the calculation into a proof.

\section{A prefix-order automaton}
\label{sec:automaton}

\subsection{General row count and the Gram statistic}
Fix $m\ge1$ and put
\[
 d=\frac{m(m+1)}2,
 \qquad \mathcal G_m=\Sym_m(\F).
\]
For $\mathcal A\subseteq\mathcal G_m$, let $N_{m,\mathcal A}(n)$ count matrices
$M\in\F^{m\times n}$ with strictly increasing rows and $MM^\top\in\mathcal A$.
The special set for our problem is
\[
 \Astar=\{G\in\mathcal G_4:
 G_{1,*}\revlex G_{2,*}\revlex G_{3,*}\revlex G_{4,*}\},
 \qquad a(n)=N_{4,\Astar}(n).
\]
Allowing only symmetric targets loses nothing, since every Gram matrix here is
symmetric. If an allowed set initially contains nonsymmetric matrices, they can
simply be discarded.

Write a matrix as a sequence of columns $v_1,\ldots,v_n\in\F^m$. Its Gram
matrix is
\begin{equation}\label{eq:gramsum}
 MM^\top=\sum_{t=1}^n v_tv_t^\top.
\end{equation}
Thus appending a column $v$ adds $vv^\top$ to the current Gram matrix.

\subsection{Remembering the first differences}
After reading a prefix of the columns, record the subset
\[
 S\subseteq\{1,\ldots,m-1\}
\]
of adjacent row pairs whose first difference has already occurred in the
correct direction. On every retained prefix, $i\in S$ means that row $i$ is
already lexicographically smaller than row $i+1$; $i\notin S$ means their
prefixes are still equal. Once resolved, a pair never needs to be compared
again.

A new column $v=(v_1,\ldots,v_m)$ is admissible from $S$ exactly when
\begin{equation}\label{eq:admissible}
 v_i\le v_{i+1}\quad\hbox{for every }i\notin S.
\end{equation}
For an admissible column, the new state is
\begin{equation}\label{eq:transition}
 \tau(S,v)=S\cup\{i:v_i=0,\ v_{i+1}=1\}.
\end{equation}
The initial state is $\varnothing$ and the accepting order state is
$S_{\rm full}=\{1,\ldots,m-1\}$. There are $2^{m-1}$ order states.

\begin{lemma}[Correctness of the order automaton]\label{lem:order}
Length-$n$ paths from $\varnothing$ to $S_{\rm full}$, with columns as transition
labels and with admissibility \eqref{eq:admissible}, are in bijection with binary
$m\times n$ matrices having strictly increasing rows.
\end{lemma}
\begin{proof}
At the empty prefix all adjacent pairs are equal. Suppose the description of a
state is correct for a retained prefix. For an unresolved pair, a new pair of
bits $(1,0)$ makes the first difference wrong and must be rejected; $(0,0)$ or
$(1,1)$ preserves equality; $(0,1)$ resolves it correctly. For a resolved pair,
any later bits leave its first difference unchanged. This proves the invariant
and the transition rule by induction.

At the end, every adjacent pair must be resolved to obtain strict order.
Conversely, any strictly increasing matrix has a first difference of type
$(0,1)$ for each adjacent pair and therefore follows exactly one accepted path.
Adjacent strict inequalities are equivalent to the whole increasing row chain.
\end{proof}

\subsection{An exact integer dynamic program}
Let $D_n(S,G)$ count retained prefixes of length $n$ with order state $S$ and
Gram matrix $G$. Start with
\begin{equation}\label{eq:dpinitial}
 D_0(\varnothing,0)=1,
 \qquad D_0(S,G)=0\quad\hbox{otherwise}.
\end{equation}
For each admissible $v$, perform the update
\begin{equation}\label{eq:dp}
 D_{n+1}\bigl(\tau(S,v),G+vv^\top\bigr)
 \mathrel{+}=D_n(S,G).
\end{equation}
Then
\begin{equation}\label{eq:dpend}
 N_{m,\mathcal A}(n)=\sum_{G\in\mathcal A}D_n(S_{\rm full},G).
\end{equation}
These equations, together with \Cref{lem:order} and \eqref{eq:gramsum}, prove
correctness of the finite enumeration used later. Different columns that lead
to the same next state must be counted separately; the update does so.

An important independent total is
\begin{equation}\label{eq:unrestricted}
 \sum_{G\in\mathcal G_m}D_n(S_{\rm full},G)=\binom{2^n}{m}.
\end{equation}
Indeed, a strictly increasing list of $m$ binary words of length $n$ is just an
$m$-element subset of the $2^n$ possible words, written in its unique increasing
order. The supplied implementation checks this equality at every computed
length.

\section{Fourier decomposition and quadratic Walsh sums}
\label{sec:fourier}

\subsection{Characters of the symmetric-matrix group}
Identify $\mathcal G_m$ with $\F^d$ using its upper-triangular entries. For
$Q=(Q_{ij})_{i\le j}\in\F^d$ and $G\in\mathcal G_m$, define
\begin{equation}\label{eq:pairing}
 \ip{Q}{G}=\sum_{i\le j}Q_{ij}G_{ij}\pmod2,
 \qquad \chi_Q(G)=(-1)^{\ip{Q}{G}}.
\end{equation}
This is an upper-triangle pairing, \emph{not} the ordinary trace pairing of two
symmetric matrices in characteristic two. Using a trace pairing would count
off-diagonal entries twice and destroy the required character group.

Let
\[
 \widehat{\1_{\mathcal A}}(Q)=\sum_{G\in\mathcal A}\chi_Q(G).
\]
The elementary orthogonality identity
\[
 \sum_{Q\in\F^d}(-1)^{\ip{Q}{H}}
 =\begin{cases}2^d,&H=0,\\0,&H\ne0\end{cases}
\]
implies
\begin{equation}\label{eq:inversion}
 \1_{\mathcal A}(G)
 =2^{-d}\sum_{Q\in\F^d}\widehat{\1_{\mathcal A}}(Q)\chi_Q(G).
\end{equation}
For completeness, when $H\ne0$, choose a coordinate where it is 1 and pair each
$Q$ with the vector obtained by toggling that coordinate; the two signs cancel.
Applying this to $H=G+G'$ proves the inversion formula.

\subsection{Small weighted transition matrices}
For $Q\in\F^d$, introduce the quadratic Boolean function
\begin{equation}\label{eq:qQ}
 q_Q(v)=\ip{Q}{vv^\top}
 =\sum_iQ_{ii}v_i+\sum_{i<j}Q_{ij}v_iv_j.
\end{equation}
The diagonal terms are linear because $v_i^2=v_i$ on $\F$.
Define the $2^{m-1}\times2^{m-1}$ integer matrix $T_Q$ by
\begin{equation}\label{eq:TQ}
 (T_Q)_{S,S'}=
 \sum_{\substack{v\text{ admissible from }S\\\tau(S,v)=S'}}(-1)^{q_Q(v)}.
\end{equation}
Vectors $\e_S$ are standard coordinate columns. The weighted number of
accepted paths is
\[
 b_Q(n)=\e_{\varnothing}^{\top}T_Q^n\e_{S_{\rm full}}.
\]
Since $\chi_Q$ turns the Gram sum into a product of column weights,
\eqref{eq:inversion} gives
\begin{equation}\label{eq:Fouriercount}
 N_{m,\mathcal A}(n)=2^{-d}
 \sum_{Q\in\F^d}\widehat{\1_{\mathcal A}}(Q)\,
 \e_{\varnothing}^{\top}T_Q^n\e_{S_{\rm full}}.
\end{equation}
For four rows, this replaces one 8192-state system by 1024 integer matrices of
size eight. More importantly, all these small matrices have tightly controlled
spectra and the same short filtration.

\subsection{A self-contained Walsh-sum lemma}
\begin{lemma}[Possible quadratic sums]\label{lem:walsh}
Let $q:\F^k\to\F$ be a polynomial function of degree at most two with $q(0)=0$,
and put
\[
 W(q)=\sum_{v\in\F^k}(-1)^{q(v)}.
\]
Then
\begin{equation}\label{eq:Walshset}
 W(q)\in\Lambda_k:=\{0,2^k\}\cup
 \{\,\pm2^{k-s}:1\le s\le\lfloor k/2\rfloor\,\}.
\end{equation}
\end{lemma}
\begin{proof}
The polarization
\[
 B(u,v)=q(u+v)+q(u)+q(v)
\]
is an alternating bilinear form. Its rank $r$ is even. One elementary proof is
by induction: if $B$ is nonzero, choose $u,v$ with $B(u,v)=1$; their span is a
nondegenerate two-dimensional summand, and its orthogonal complement reduces
the rank by two. If $B=0$, its rank is zero.

Let $R=\rad B$, so $\dim R=k-r$. Squaring the sum and substituting $w=v+u$
gives
\begin{align*}
 W(q)^2
 &=\sum_{u,v\in\F^k}(-1)^{q(v)+q(v+u)}\\
 &=\sum_{u\in\F^k}(-1)^{q(u)}
       \sum_{v\in\F^k}(-1)^{B(v,u)}
 =2^k\sum_{u\in R}(-1)^{q(u)}.
\end{align*}
On $R$ the function $q$ is linear, since its polarization vanishes there.
The last sum is either zero or $2^{k-r}$. Hence $W(q)=0$ or
$W(q)=\pm2^{k-r/2}$.

There is one sign restriction to retain. If $r=0$, then $q$ is linear on all
of $\F^k$. Its sum is zero when it is nonzero and $2^k$ when it is zero; it
cannot be $-2^k$. For positive even $r=2s$, the remaining possibilities are
exactly those permitted by \eqref{eq:Walshset}. The argument only needs this
inclusion, not that every listed value occurs in every application.
\end{proof}

\subsection{Why the diagonal entries are Walsh sums}
A transition can stay at the same order state $S$ only when
\[
 v_i=v_{i+1}\qquad(i\notin S).
\]
These equations say that $v$ is constant on each contiguous block separated by
the boundaries in $S$. There are $k=|S|+1$ such blocks. Assigning a bit to each
block parametrizes the allowed staying columns by $\F^k$. Restricting $q_Q$ to
this subspace preserves its degree bound and its zero constant term. Therefore
\begin{equation}\label{eq:diagonal}
 (T_Q)_{S,S}\in\Lambda_{|S|+1}.
\end{equation}
Also a nonzero transition only moves from $S$ to a superset $S'$. In particular,
it never moves to a different state on the same cardinality level. This second
fact controls multiplicities; the diagonal values alone would not suffice.

\section{An explicit universal recurrence}
\label{sec:universal}

\subsection{An annihilator from the level filtration}
\begin{lemma}[Filtered annihilation]\label{lem:filtration}
Suppose $T$ preserves a filtration
$0=V_0\subseteq V_1\subseteq\cdots\subseteq V_m=V$.
If the induced operator on $V_k/V_{k-1}$ is annihilated by a polynomial $p_k$,
then $\prod_{k=1}^mp_k(T)=0$.
\end{lemma}
\begin{proof}
The assumption is exactly $p_k(T)V_k\subseteq V_{k-1}$. Apply $p_m(T)$ first,
then $p_{m-1}(T)$, down to $p_1(T)$. Their product maps $V_m$ into $V_0=0$.
All factors are polynomials in the same operator, so they commute; their written
order is immaterial.
\end{proof}

Order the coordinate basis by $|S|$. Regard $T_Q$ as acting on column vectors,
and let $V_k$ be the span of $\e_S$ with $|S|\le k-1$. Because
$(T_Q)_{S,S'}$ can be nonzero only when $S\subseteq S'$, these subspaces are
invariant. On the quotient $V_k/V_{k-1}$ the matrix is diagonal. Thus
\eqref{eq:diagonal} and \Cref{lem:filtration} apply with
\begin{equation}\label{eq:pk}
 p_k(x)=\prod_{\lambda\in\Lambda_k}(x-\lambda)
       =x(x-2^k)\prod_{s=1}^{\lfloor k/2\rfloor}
         (x-2^{k-s})(x+2^{k-s}).
\end{equation}
The polynomial has only one copy of each allowed diagonal value at that level.
We do not multiply a characteristic polynomial for every state.

\begin{theorem}[Universal spectral recurrence]\label{thm:universal}
For $m\ge1$, set $r_j=\min(j,m-j)$ and define
\begin{align}
 P_m(x)&=\prod_{j=1}^m(x-2^j)^{r_j+1}(x+2^j)^{r_j},
 \label{eq:Pm}\\
 U_m(x)&=x^mP_m(x).
 \label{eq:Um}
\end{align}
For every $\mathcal A\subseteq\mathcal G_m$,
\begin{equation}\label{eq:universalE}
 U_m(E)N_{m,\mathcal A}(n)=0\quad(n\ge0),
 \qquad
 P_m(E)N_{m,\mathcal A}(n)=0\quad(n\ge m).
\end{equation}
Moreover,
\begin{equation}\label{eq:degrees}
 \deg P_m=m+2\lfloor m^2/4\rfloor,
 \qquad \deg U_m=2m+2\lfloor m^2/4\rfloor.
\end{equation}
\end{theorem}
\begin{proof}
Equations \eqref{eq:pk} and the filtration lemma show that
$\prod_{k=1}^mp_k(T_Q)=0$ for every $Q$. We identify this product explicitly.
Zero occurs once at each level, producing $x^m$.

The positive value $2^j$ occurs at level $k=j$ through the rank-zero term,
and at levels
\[
 j+1\le k\le\min(2j,m)
\]
through a rank $2(k-j)$ quadratic sum. Its total multiplicity in the product
is $1+\min(j,m-j)$. The negative value $-2^j$ occurs at the same latter levels
but not at $k=j$, so its multiplicity is $\min(j,m-j)$.
Consequently $\prod_kp_k=U_m$ and $U_m(T_Q)=0$.

For any polynomial $R$, the action of $R(E)$ on
$\e_{\varnothing}^{\top}T_Q^n\e_{S_{\rm full}}$ is
\[
 \e_{\varnothing}^{\top}T_Q^nR(T_Q)\e_{S_{\rm full}}.
\]
Thus $U_m(E)$ annihilates every Fourier summand in \eqref{eq:Fouriercount},
and hence the count itself. Since $U_m(E)=E^mP_m(E)$, the second assertion in
\eqref{eq:universalE} follows with the stated index range.
Finally,
\[
 \deg P_m=\sum_{j=1}^m(2\min(j,m-j)+1),
 \qquad \sum_{j=1}^m\min(j,m-j)=\lfloor m^2/4\rfloor,
\]
which yields \eqref{eq:degrees}.
\end{proof}

\begin{corollary}[Form of every tail]\label{cor:tail}
For $n\ge m$, there are polynomials $A_j,B_j\in\Q[n]$ such that
\begin{equation}\label{eq:generalform}
 N_{m,\mathcal A}(n)
 =\sum_{j=1}^m\bigl(A_j(n)2^{jn}+B_j(n)(-2^j)^n\bigr),
\end{equation}
with
\[
 \deg A_j\le\min(j,m-j),\qquad
 \deg B_j\le\min(j,m-j)-1.
\]
A negative degree bound means the polynomial is zero. The generating function
is rational. A candidate of the form \eqref{eq:generalform} is certified by
checking exactly $\deg P_m$ consecutive values beginning at $n=m$.
\end{corollary}
\begin{proof}
For $\lambda\ne0$,
\[
 (E-\lambda)\bigl(\lambda^np(n)\bigr)
 =\lambda^{n+1}\bigl(p(n+1)-p(n)\bigr).
\]
Repeated application lowers the polynomial degree. The sequences
$n^h\lambda^n$ for distinct nonzero $\lambda$ and the prescribed degree ranges
are linearly independent and give $\deg P_m$ solutions of the monic recurrence.
Independence follows, for example, by applying the factors belonging to all
roots except one and then taking finite differences at the remaining root.
A recurrence of order $\deg P_m$ is uniquely determined by that many initial
values, so these sequences span its solution space. The coefficients are
rational because the corresponding linear system has rational entries and a
unique solution. Their generating functions are rational; changing a finite
prefix only adds a polynomial. The same uniqueness proves the certification
claim.
\end{proof}

\subsection{The four-row specialization}
The level spectra are particularly small:
\begin{center}
\begin{tabular}{ccl}
\toprule
Level $k$ & Number of states & Permitted staying sums $\Lambda_k$\\
\midrule
1 & 1 & $0,2$\\
2 & 3 & $0,-2,2,4$\\
3 & 3 & $0,-4,4,8$\\
4 & 1 & $0,-8,-4,4,8,16$\\
\bottomrule
\end{tabular}
\end{center}
Multiplying the four level polynomials gives $U_4=x^4P_4$, with $P_4$ exactly
as in \eqref{eq:P4}. Therefore \eqref{eq:keyrec} is proved for the target count
and, in fact, for every four-row Gram acceptance condition. The order 12 is a
consequence of the combinatorial construction and the quadratic-sum lemma;
it is not selected because twelve values happened to fit a formula.

\section{The finite certificate and proof of the formula}
\label{sec:certificate}

\subsection{Encoding the finite calculation}
For reproducibility, the bit-mask implementation orders the ten Gram
coordinates as
\begin{equation}\label{eq:coordinateorder}
 (11),(12),(13),(14),(22),(23),(24),(33),(34),(44).
\end{equation}
The first coordinate is the least significant bit of the Gram mask. State bit
$i-1$ says that adjacent rows $i,i+1$ have been resolved. A column mask uses its
least significant bit for its top entry. These choices affect only the internal
encoding, not the lexicographic convention for rows.

The acceptance test reconstructs each symmetric Gram matrix, reads each row
from left to right, and checks strict decrease. It gives exactly 48 acceptable
matrices. A useful small classification is
\begin{center}
\begin{tabular}{ccr}
\toprule
First row & Binary integer & Number of acceptable completions\\
\midrule
$1000$ & 8  & 9\\
$1100$ & 12 & 14\\
$1110$ & 14 & 18\\
$1111$ & 15 & 7\\
\midrule
Total &&48\\
\bottomrule
\end{tabular}
\end{center}
There are no other first rows. Indeed, decreasing rows imply a nonincreasing
first column, symmetry identifies that column with the first row, and an
all-zero first row cannot be greater than the second row. For each of the four
possibilities, the six remaining upper-triangular bits give just 64 completions
to test. The full list of accepted masks is included in the JSON certificate.

\begin{certificate}[Definition-based integer counts]\label{cert:counts}
The dynamic program \eqref{eq:dpinitial}--\eqref{eq:dpend} gives the following
values. The auxiliary column $z(n)$ counts the same increasing-row matrices
with $MM^\top=0$; it is used only for the sharpness result in
\Cref{sec:sharpness}.
\begin{center}
\begin{tabular}{r@{\qquad}r@{\qquad}r}
\toprule
$n$ & $a(n)$ & $z(n)$\\
\midrule
4  & $58$                & $3$\\
5  & $1\,629$            & $15$\\
6  & $28\,924$           & $1\,020$\\
7  & $507\,052$          & $8\,820$\\
8  & $8\,211\,776$      & $202\,608$\\
9  & $133\,693\,904$    & $2\,656\,080$\\
10 & $2\,140\,571\,200$ & $46\,626\,240$\\
11 & $34\,361\,115\,072$ & $706\,770\,240$\\
12 & $549\,587\,348\,992$ & $11\,575\,711\,488$\\
13 & $8\,798\,356\,254\,976$ & $182\,670\,808\,320$\\
14 & $140\,744\,002\,571\,264$ & $2\,939\,896\,028\,160$\\
15 & $2\,252\,082\,614\,856\,704$ & $46\,875\,280\,757\,760$\\
\bottomrule
\end{tabular}
\end{center}
At every displayed index, exact rational evaluation of the right side of
\eqref{eq:formula} equals $a(n)$. The same enumeration gives zero at
$n=0,1,2,3$ for both columns.
\end{certificate}

The complete small program in \Cref{app:code} calculates the $a(n)$ column
without importing any tabulated sequence values. It uses tuple states and
tuple Gram coordinates. The separate program \texttt{verify.py} uses bit masks,
also calculates $z(n)$, and performs the additional audits in
\Cref{sec:computation}. Thus the certificate is an explicit finite arithmetic
claim, not an appeal to the authority of a pre-existing data table.

\subsection{Why the twelve checks finish the proof}
Let $F(n)$ denote the expression on the right side of \eqref{eq:formula}, and
put
\begin{equation}\label{eq:Q}
 Q(x)=(x-16)(x-8)^2(x+8)(x-4)^3(x+4)^2(x-2)^2.
\end{equation}
The finite-difference identity in the proof of \Cref{cor:tail} shows directly
that $Q(E)F=0$ for all nonnegative indices. As $P_4(x)=(x+2)Q(x)$, also
$P_4(E)F=0$.

\begin{proof}[Proof of \Cref{thm:main}]
By \Cref{thm:universal}, $P_4(E)a(n)=0$ for $n\ge4$. The preceding observation
shows that the difference $h(n)=a(n)-F(n)$ satisfies that same order-12 monic
recurrence on this range. \Cref{cert:counts} gives
$h(4)=h(5)=\cdots=h(15)=0$.

Write $P_4(x)=x^{12}+\sum_{j=0}^{11}p_jx^j$. The recurrence at $n=4$ gives
$h(16)=-\sum_{j=0}^{11}p_jh(4+j)=0$. Inductively, if $h$ vanishes through
$n+11$, the recurrence at $n$ gives $h(n+12)=0$. Hence $h(n)=0$ for every
$n\ge4$. The initial zeros come from the finite certificate, or directly from
the same exact enumeration for these four small lengths.
\end{proof}

\begin{remark}[The transient is real]
The lower limit $n\ge4$ must not be dropped. The extrapolated expression gives
\[
 F(0)=\frac{103}{4096},\qquad F(1)=-\frac{71}{1024},\qquad
 F(2)=-\frac1{64},\qquad F(3)=-\frac1{16},
\]
whereas the counting sequence is zero at all four indices. The zero root in
the universal annihilator accounts for such finite transients.
\end{remark}

\section{Generating functions, minimal recurrences, and sharpness}
\label{sec:gf}

\subsection{The exact rational generating function}
Define
\begin{equation}\label{eq:D}
 D(z)=(1-16z)(1-8z)^2(1+8z)(1-4z)^3(1+4z)^2(1-2z)^2
\end{equation}
and
\begin{align}
 B(z)={}&58-227z-9052z^2+68976z^3-16960z^4 \notag\\
 &-1025408z^5+3055104z^6-3796992z^7 \notag\\
 &+2834432z^8-2490368z^9+2097152z^{10}.
 \label{eq:B}
\end{align}

\begin{corollary}[Rational generating function]\label{cor:gf}
As an identity of formal power series,
\begin{equation}\label{eq:GF}
 \sum_{n\ge0}a(n)z^n=z^4\frac{B(z)}{D(z)}.
\end{equation}
Equivalently, $\sum_{k\ge0}a(k+4)z^k=B(z)/D(z)$.
\end{corollary}
\begin{proof}
By \Cref{thm:main}, the shifted sequence $b(k)=a(k+4)$ is annihilated by
$Q(E)$ for every $k\ge0$. The reciprocal characteristic polynomial is
$D(z)=z^{11}Q(1/z)$. If $D(z)=\sum_{j=0}^{11}d_jz^j$, the coefficients of
$D(z)\sum_{k\ge0}b(k)z^k$ vanish from degree 11 onward. For $0\le k\le10$ they
are the finite convolutions
\[
 \sum_{j=0}^kd_j a(k-j+4).
\]
Using \Cref{cert:counts} gives exactly \eqref{eq:B}. The four initial zero terms
then supply the factor $z^4$.
\end{proof}

For direct implementation the characteristic polynomial is
\begin{align}
 Q(x)={}&x^{11}-32x^{10}+244x^9+1552x^8-27392x^7\notag\\
 &+70144x^6+484352x^5-2781184x^4+1212416x^3\notag\\
 &+20185088x^2-48234496x+33554432.
 \label{eq:Qexpanded}
\end{align}
Thus $\sum_{j=0}^{11}q_ja(n+j)=0$ for $n\ge4$, where $q_j$ is the coefficient
of $x^j$ in \eqref{eq:Qexpanded}. The range is important: this order-11
recurrence fails at each of $n=0,1,2,3$.

\subsection{Minimality for the selected sequence}
\begin{proposition}\label{prop:minimal}
The minimal monic constant-coefficient annihilator of the eventual tail of
$a(n)$ is $Q(x)$, of degree 11. If a recurrence is required to hold at every
$n\ge0$, its minimal monic annihilator is $x^4Q(x)$, of degree 15.
\end{proposition}
\begin{proof}
In \eqref{eq:formula}, the six distinct bases and the exact polynomial degrees
are
\[
\begin{array}{c|rrrrrr}
\lambda&16&8&-8&4&-4&2\\\hline
\deg p_\lambda&0&1&0&2&1&1.
\end{array}
\]
Each leading coefficient is nonzero. The independence and finite-difference
argument from \Cref{cor:tail} imply that a polynomial annihilates
$p_\lambda(n)\lambda^n$ only if it contains $(x-\lambda)$ to at least
$1+\deg p_\lambda$. Distinct roots cannot cancel one another's components.
Therefore the minimal eventual annihilator is precisely \eqref{eq:Q}.

For the global assertion, write $a=F+\delta$, where $\delta$ is supported on
$0,1,2,3$ and $\delta(3)=1/16\ne0$. The forward shift on sequences supported in
this range is nilpotent, and $\delta$ has nilpotency index exactly four. Any
global annihilator must be divisible by $Q$, because its action on the
exponential-polynomial tail must vanish. Since $Q(0)\ne0$, the operator $Q(E)$
is invertible on this finite-support space and preserves that nilpotency
index. Consequently a further factor $x^4$ is necessary and sufficient to
kill the transient. This proves the second assertion.
\end{proof}

There is no conflict between an eventual recurrence order of 11 here and the
order-10, degree-one polynomial-coefficient guess reported in the 2023 paper:
constant-coefficient order and polynomial-coefficient order are different
notions. The minimality assertion above concerns constant coefficients only.

\subsection{The four-row universal bound is sharp}
\label{sec:sharpness}
It is natural to ask whether the factor $x+2$ in $P_4$ was only an artifact of
the argument: it disappears from the answer for $\Astar$. It is not an artifact
for arbitrary acceptance sets.

\begin{proposition}[Sharp universal order for four rows]\label{prop:sharp}
Among all polynomials that eventually annihilate $N_{4,\mathcal A}$ for
\emph{every} $\mathcal A\subseteq\mathcal G_4$, the minimal monic one is $P_4$.
Its degree is 12. If annihilation must hold from $n=0$, the minimal common
monic polynomial is $x^4P_4$, of degree 16.
\end{proposition}
\begin{proof}
Theorem \ref{thm:universal} supplies both upper bounds. A common eventual
annihilator must be divisible by $Q$, by \Cref{prop:minimal}. To see the missing
factor, use the target $\mathcal A=\{0\}$ and write
$z(n)=N_{4,\{0\}}(n)$.

The universal theorem says that on $n\ge4$, the only component of $z(n)$ not
annihilated by $Q(E)$ is a constant multiple $c(-2)^n$, because the multiplicity
of $-2$ in $P_4$ is one. The exact counts in \Cref{cert:counts}, substituted
into \eqref{eq:Qexpanded}, give
\[
 Q(E)z(4)=136857600,
 \qquad Q(-2)=149299200.
\]
It follows that
\begin{equation}\label{eq:minus2}
 c=\frac{136857600}{(-2)^4\cdot149299200}=\frac{11}{192}\ne0.
\end{equation}
Hence $x+2$ is also necessary. Its product with $Q$ is $P_4$, proving eventual
sharpness.

For a common annihilator valid from $n=0$, the selected sequence $a$ already
forces $x^4Q$ by \Cref{prop:minimal}; the zero-Gram target forces $x+2$.
Their least common multiple is $x^4P_4$, completing the proof.
\end{proof}

The degrees 11, 12, 15, and 16 therefore answer four different questions, not
four inconsistent descriptions of the same recurrence.

\section{Equidistribution and exact asymptotic information}
\label{sec:asymptotics}

\subsection{A general uniformity theorem}
\begin{theorem}[Ordered Gram equidistribution]\label{thm:uniform}
For each fixed $m\ge2$ and every $\mathcal A\subseteq\mathcal G_m$,
\begin{equation}\label{eq:uniformcount}
 N_{m,\mathcal A}(n)
 =\frac{|\mathcal A|}{2^{m(m+1)/2}}\binom{2^n}{m}
  +O_m\!\left(n\,2^{(m-1)n}\right)
 \quad(n\to\infty).
\end{equation}
In particular, if an increasing-row matrix is chosen uniformly, its Gram
matrix tends to the uniform distribution on all symmetric $m\times m$ binary
matrices, and
\begin{equation}\label{eq:uniformprob}
 \Pr(MM^\top\in\mathcal A)
 =\frac{|\mathcal A|}{2^{m(m+1)/2}}+O_m(n2^{-n}).
\end{equation}
The implied constants may be chosen uniformly over $\mathcal A$ for fixed $m$.
\end{theorem}
\begin{proof}
The zero character in \eqref{eq:Fouriercount} contributes exactly
$2^{-d}|\mathcal A|\binom{2^n}{m}$, by \eqref{eq:unrestricted}.
Consider a nonzero character $Q$.
At every level $k<m$, a diagonal sum has absolute value at most $2^k\le2^{m-1}$.
At level $m$, the function $q_Q$ in \eqref{eq:qQ} is nonzero: its values on
$e_i$ recover $Q_{ii}$, and its values on $e_i+e_j$, together with those
single-coordinate values, recover $Q_{ij}$. Thus $q_Q$ cannot be identically
zero unless $Q=0$.

If the polarization of this nonzero $q_Q$ has rank zero, its Walsh sum is zero.
Otherwise its rank is at least two, and \Cref{lem:walsh} bounds the sum by
$2^{m-1}$ in absolute value. Hence no nonzero character has the eigenvalue
$2^m$.

There is also a sharp enough multiplicity bound. The value $+2^{m-1}$ can
occur only at levels $m-1$ and $m$, so its polynomial prefactor has degree at
most one. The value $-2^{m-1}$ can occur only at level $m$, so its prefactor is
constant. Every other nonzero eigenvalue has smaller absolute value and a
bounded polynomial prefactor, and a zero-eigenvalue component has finite
support. Each nonzero Fourier summand is consequently
$O_m(n2^{(m-1)n})$. There are finitely many characters, with
$|\widehat{\1_{\mathcal A}}(Q)|\le2^d$, so summing preserves this estimate
uniformly over the target set. Division by
$\binom{2^n}{m}\sim2^{mn}/m!$ proves \eqref{eq:uniformprob}.
\end{proof}

The sorting of rows is relevant here: conditional Gram uniformity should not
simply be assumed from an unconstrained random-matrix model. The proof shows
why sorting affects only exponentially smaller terms.

\subsection{Leading terms for A181280}
Because $|\Astar|=48$ and $2^d=1024$, \Cref{thm:uniform} predicts the limiting
conditional probability $48/1024=3/64$. The exact formula gives more:
\begin{align}
 a(n)={}&\frac{16^n}{512}
 +\left(\frac{n}{512}-\frac{3473}{147456}
              -\frac{113(-1)^n}{49152}\right)8^n
 +O(n^2 4^n),
 \label{eq:asympa}\\
 \frac{a(n)}{\binom{2^n}{4}}
 ={}&\frac3{64}
 +\left(\frac{3n}{64}-\frac{1745}{6144}
              -\frac{113(-1)^n}{2048}\right)2^{-n}
 +O(n^2 4^{-n}).
 \label{eq:asympprob}
\end{align}
Thus $a(n)\sim16^n/512$. Among \emph{all} binary $4\times n$ matrices, the
probability of satisfying both ordering conditions tends to $1/512$; among
those with already increasing rows it tends to $3/64$. These are different
sample spaces and should not be conflated.

The original count has a finite exact exponential-polynomial expansion:
\eqref{eq:formula} terminates after absolute exponential scales $16^n,8^n,4^n$
and $2^n$, with parity effects at the middle two scales. There is no unproved
asymptotic remainder hidden in that identity.

\subsection{All coefficients of the conditional-probability expansion}
A compact exact rational form gives an infinite expansion for the conditional
probability as well. Put $x=2^{-n}$ and $\varepsilon=(-1)^n$, and define
\begin{align*}
 \nu_0&=\frac3{64},\\
 \nu_1(n,\varepsilon)&=\frac{288n-3473}{6144}
                         -\frac{113\varepsilon}{2048},\\
 \nu_2(n,\varepsilon)&=\frac{6n^2-219n+820-\varepsilon(13n-164)}{256},\\
 \nu_3(n,\varepsilon)&=-\frac{3n+32}{12}.
\end{align*}
Since $\binom{2^n}{4}=2^{4n}(1-x)(1-2x)(1-3x)/24$, the exact identity is
\begin{equation}\label{eq:probrational}
 \frac{a(n)}{\binom{2^n}{4}}
 =\frac{\nu_0+\nu_1x+\nu_2x^2+\nu_3x^3}{(1-x)(1-2x)(1-3x)}
 \qquad(n\ge4).
\end{equation}
Let $h_t=0$ for $t<0$ and, for $t\ge0$, let
\begin{equation}\label{eq:hcoef}
 h_t=\frac{1-2^{t+3}+3^{t+2}}2.
\end{equation}
Partial fractions give
\[
 \frac1{(1-x)(1-2x)(1-3x)}
 =\frac{1/2}{1-x}-\frac4{1-2x}+\frac{9/2}{1-3x}
 =\sum_{t\ge0}h_tx^t.
\]
It follows that the coefficient polynomials in the full expansion are
\begin{equation}\label{eq:allcoeffs}
 C_k(n,\varepsilon)=\sum_{j=0}^3\nu_j(n,\varepsilon)h_{k-j},
 \qquad
 \frac{a(n)}{\binom{2^n}{4}}
 =\sum_{k\ge0}C_k(n,(-1)^n)2^{-kn}.
\end{equation}
For each fixed $n\ge4$ the latter series converges absolutely. For a fixed
truncation order $K$, its tail is $O_K(n^2 2^{-(K+1)n})$ as $n\to\infty$.
Indeed $h_t=O(3^t)$, all $\nu_j$ are $O(n^2)$, and $3\cdot2^{-n}<1$ on this
range. This supplies every coefficient without further recurrence guessing.

\section{Computation, reproducibility, and audit boundaries}
\label{sec:computation}

\subsection{What was executed}
The supplied standard-library program \texttt{verify.py} was executed using
Python 3.13.5. Its recorded run completed in approximately 4.15 seconds on the
working machine; timing is an observation, not a performance guarantee.
The output is saved in \texttt{results/verification.txt}.

The executed checks comprise exact dynamic programming through $n=100$;
comparison with \eqref{eq:formula} throughout $4\le n\le100$; the unrestricted
count \eqref{eq:unrestricted}; all 1024 weighted $8\times8$ identities
$U_4(T_Q)=0$; Fourier inversion through $n=15$; and the displayed generating
function and recurrence coefficients. Every nonzero-character matrix also
passes the spectral-radius bound eight directly on its diagonal.

A separate brute-force enumeration chooses all four-element subsets of binary
row words for $0\le n\le6$ and computes every Gram row directly using parity
of bit intersections. At $n=6$ this checks $\binom{64}{4}=635376$ row sets. It
uses neither order-state transitions nor upper-triangular Gram masks. Both the
A181280 count and the zero-Gram count agree with the dynamic program.
For $m=1,2,3$, the program additionally checks the universal recurrence through
$n=20$ for \emph{every} individual Gram target.

Finally, \texttt{minimal\_certificate.py}, reproduced in \Cref{app:code}, was
executed independently. It represents states and Gram entries by tuples rather
than masks and reproduces all twelve decisive values. The scripts use exact
integers and rational arithmetic only. They perform no network requests and
read no sequence data as input.

\subsection{What makes this a finite proof rather than a large test}
The mathematical implication is
\[
 \text{definition}
 \Longrightarrow U_4(E)a=0
 \Longrightarrow P_4(E)a=0\text{ on }n\ge4,
\]
followed by uniqueness of solutions of a monic order-12 recurrence and a finite
arithmetic certificate. The extra computations through $n=100$ are useful
cross-checks, but the all-$n$ conclusion does not depend on assuming that a
pattern observed through 100 continues.

The computer-assisted dependency is transparent: the arithmetic in
\Cref{cert:counts} and in the small acceptance enumeration must be correct.
The transition invariant, Fourier identity, quadratic-sum argument, recurrence
bound, and finite-to-infinite induction are conventional proofs given in full.
The two state encodings and the independent brute force reduce implementation
risk; they do not replace mathematical review. No Lean formalization or
external peer review is claimed.

\subsection{Complexity and practical use}
The direct state space has $2^{d+m-1}$ entries and at most $2^m$ column choices
per state. A straightforward implementation through length $N$ uses
\[
 O\!\left(N2^{d+2m-1}\right)
\]
integer additions and two layers of $O(2^{d+m-1})$ integer registers, apart from
optional stored history. Counts have $O(mN)$ bits. Thus for fixed $m$ the
arithmetic-operation count is linear in $N$, while a naive bit-operation bound
for the additions is quadratic in $N$. The full verifier retains endpoint
Gram distributions at all lengths for its audits, so its stored history uses
$O(N2^d)$ integer registers in addition to the working layers.

For the selected sequence, once the formula is established, either
\eqref{eq:formula} or the order-11 recurrence is preferable for computing very
large indices. An isolated value can be evaluated with integer exponentiation
and a fixed amount of rational arithmetic. This does not imply polylogarithmic
bit complexity: the answer itself has $\Theta(n)$ bits.

The reusable enumerator explicitly limits its implementation to $1\le m\le5$
to avoid accidental exponential resource consumption. This is a software guard,
not a limitation of \Cref{thm:universal}. The present exhaustive Fourier audit
was performed for $m=4$, not for all $m$.

\section{Conclusions and remaining research directions}
\label{sec:conclusion}
The explicit A181280 identity has a complete proof here, with a small,
reproducible computer-assisted arithmetic component. The decisive step is the
universal spectral bound, not a more aggressive fit of the existing data.
The factor $(x+2)$ vanishes for the decreasing-Gram condition but is necessary
for other targets, and the four exceptional initial indices are explained by
a genuine nilpotent transient. The rational generating function, minimal
recurrences, and asymptotic probability all follow from the same calculation.

For general $m$, the theorem gives a uniform quadratic-size recurrence bound
despite an exponentially large finite-state realization. What remains to be
determined by further work is when that universal bound is sharp for other
row counts, which target conditions cancel particular spectral components,
and whether explicit structural formulas can replace the remaining finite
arithmetic for broad families of target sets. These are directions suggested
by the present results, not additional open problems asserted to have been
settled here.

The status distinction remains important. C-finiteness was already proved in
the retrieved 2025 slides, and the explicit formula was already conjectured in
2023. The present package supplies its proof and the additional structural
results developed above. Establishing priority would require a fuller
literature and author check than the available-source audit carried out here.

\appendix
\section{The recurrence certificate in machine-readable form}
\label{app:recurrences}
The file \texttt{results/certificate.json} records the twelve initial values,
the 48 accepted Gram masks in the coordinate order \eqref{eq:coordinateorder},
all recurrence and generating-function coefficient arrays, the independent
brute-force results, and the Fourier audit summary. Polynomial arrays are
always listed in increasing powers.

For a convenient manual comparison, the universal eventual polynomial is
\begin{align*}
 P_4(x)={}&x^{12}-30x^{11}+180x^{10}+2040x^9-24288x^8
                +15360x^7\\
          &+624640x^6-1812480x^5-4349952x^4+22609920x^3\\
          &-7864320x^2-62914560x+67108864.
\end{align*}
The four nonzero early residuals of the smaller recurrence are
\[
 \bigl(Q(E)a(0),Q(E)a(1),Q(E)a(2),Q(E)a(3)\bigr)
 =(-3796992,2834432,-2490368,2097152).
\]
They provide an especially simple check that an implementation has not silently
shifted the recurrence's valid range.

The observed level spectra in the exhaustive Fourier audit were exactly
$\Lambda_1,\Lambda_2,\Lambda_3,\Lambda_4$ as listed in
\Cref{sec:universal}. This observation is supplementary: the proof only needs
containment in those sets, which follows analytically from
\Cref{lem:walsh}.

\clearpage
\section{A complete small arithmetic verifier}
\label{app:code}
The following is the complete file \texttt{minimal\_certificate.py}. Its
ordinary list comparisons implement left-to-right lexicographic order.
It proves the finite arithmetic portion of \Cref{thm:main}; the universal
recurrence theorem supplies the all-$n$ implication.
\lstinputlisting[basicstyle=\ttfamily\fontsize{8}{9.1}\selectfont]{minimal_certificate.py}

\section{Package contents and reproduction}
The main article is \texttt{article.tex}, compiled to \texttt{article.pdf}.
Both Python programs need only Python 3.10 or later. The \texttt{results/} directory contains the generated
JSON certificate, CSV counts through $n=100$, and executed audit output.
\texttt{STATUS.md} records the source-status distinctions; \texttt{README.md}
contains usage instructions; \texttt{build.sh} rebuilds this PDF with pdfLaTeX.

From the package root, the principal commands are
\begin{lstlisting}[language=bash]
python3 minimal_certificate.py
python3 verify.py --out results
sh build.sh
\end{lstlisting}
The article does not depend on a symbolic-algebra package. The TeX source uses
standard TeX Live packages, including \texttt{newtx}, \texttt{amsmath},
\texttt{tcolorbox}, \texttt{listings}, and \texttt{hyperref}. No font files or
third-party source documents are included in the archive.

\begin{thebibliography}{9}
\bibitem{oeis}
R. H. Hardin,\
\emph{A181280: Number of $4\times n$ binary matrices with increasing rows and
mod-2 Gram matrix with decreasing rows},
The On-Line Encyclopedia of Integer Sequences.
Entry introduced 10 October 2010; formula attributed to M. Kauers and
C. Koutschan on 2 March 2023.
\url{https://oeis.org/A181280}.
Retrieved 19 September 2026.

\bibitem{kk}
M. Kauers and C. Koutschan,
\emph{Some D-Finite and Some Possibly D-Finite Sequences in the OEIS},
Journal of Integer Sequences \textbf{26} (2023), Article 23.4.5.
See Section 6.5, Conjecture 20.
\url{https://cs.uwaterloo.ca/journals/JIS/VOL26/Koutschan/kout4.html}.
Author version: arXiv:2303.02793v2, 24 April 2023,
\url{https://arxiv.org/html/2303.02793v2}.
Retrieved 19 September 2026.

\bibitem{slides}
M. Kauers,
\emph{Guessing with little data},
Rutgers Experimental Mathematics Seminar, 11 December 2025, lecture slides.
Formula on slide 30 (PDF page 104); general C-finiteness theorem on slide 39
(PDF page 137); overview check marks on slides 24 and 40.
\url{https://sites.math.rutgers.edu/~zeilberg/expmath/kauers25.pdf}.
Retrieved 19 September 2026.

\bibitem{seminar}
Rutgers Experimental Mathematics Seminar,
\emph{Archive of Speakers and Talks---2025}, entry for 11 December 2025.
\url{https://sites.math.rutgers.edu/~zeilberg/expmath/archive25.html}.
Retrieved 19 September 2026.
\end{thebibliography}
\end{document}
